\documentclass[12pt]{article}
\usepackage{fontspec}
\setmainfont{Times New Roman}
\usepackage{amsmath, amssymb, amsthm}
\usepackage{geometry}
\geometry{a4paper, margin=1in}

\title{Лекция по Вероятности и Статистика: Елементи на Комбинаториката и Случайни Блуждаения}
\author{}
\date{}

\begin{document}

\maketitle

\section{Въведение в Комбинаториката}
В тази лекция ще разгледаме някои основни елементи на комбинаториката, които са фундаментални за теорията на вероятностите. 

\subsection{Правило за умножение и пермутации}
Нека разгледаме крайно множество от $n$ елемента. Броят на възможните наредби на тези елемента (пермутации без повторение) се дава от факториела на $n$:
\[ n! = n(n-1) \dots 2 \cdot 1 \]
По дефиниция приемаме, че $0! = 1$.

\subsection{Биномни коефициенти (Комбинации)}
Ако искаме да изберем $k$ елемента от множество с $n$ елемента, като редът на избиране няма значение (комбинации без повторение), броят на възможните избори се означава с биномния коефициент:
\[ \binom{n}{k} = C_n^k = \frac{n!}{k!(n-k)!} = \frac{n(n-1)\dots(n-k+1)}{k!} \]
Тези числа се наричат още биномни коефициенти, тъй като участват в Нютоновия бином.

Основни свойства на биномните коефициенти:
1. Симетрия:
\[ \binom{n}{k} = \binom{n}{n-k} \]
2. Рекурентна връзка (Правило на Паскал):
\[ \binom{n}{k} = \binom{n-1}{k-1} + \binom{n-1}{k} \]
Тази формула позволява лесно пресмятане на биномните коефициенти и конструирането на Триъгълника на Паскал.

3. Сума на биномните коефициенти (от Нютоновия бином при $a=1, b=1$):
\[ \sum_{k=0}^{n} \binom{n}{k} = \binom{n}{0} + \binom{n}{1} + \dots + \binom{n}{n} = 2^n \]
Това число представлява броя на всички подмножества на множество с $n$ елемента.

4. Знакопроменлива сума на биномните коефициенти (от Нютоновия бином при $a=1, b=-1$):
\[ \sum_{k=0}^{n} (-1)^k \binom{n}{k} = \binom{n}{0} - \binom{n}{1} + \dots + (-1)^n \binom{n}{n} = 0 \]
Това означава, че броят на подмножествата с четен брой елементи е равен на броя на подмножествата с нечетен брой елементи.

\subsection{Нютонов бином}
Формулата за Нютоновия бином гласи:
\[ (a+b)^n = \sum_{k=0}^{n} \binom{n}{k} a^k b^{n-k} = \binom{n}{0} b^n + \binom{n}{1} ab^{n-1} + \dots + \binom{n}{n} a^n \]

\subsection{Разпределяне на частици по клетки}
Често в комбинаториката се срещат задачи за разпределяне на частици по клетки.
Ако имаме $k$ различими частици и $n$ клетки, броят на възможните разпределения е:
\[ n^k \]
Това е така, защото всяка от $k$-те частици може да бъде поставена в коя да е от $n$-те клетки напълно независимо от останалите.

\subsection{Уравнения с цели неотрицателни числа (Метод "Звезди и прегради")}
Разглеждаме уравнението:
\[ x_1 + x_2 + \dots + x_k = n \]
където $x_i$ са цели неотрицателни числа ($x_i \geq 0$ за всяко $i=1,\dots,k$) и $n \geq 0$ е цяло число.
Търсим броя на решенията на това уравнение. Можем да си представим, че имаме $n$ неразличими обекта (единици или "звезди"), които искаме да разпределим в $k$ клетки. За да разделим обектите на $k$ групи, ни трябват $k-1$ "прегради" (или "черти"). 

Общият брой на позициите за звездите и преградите е $n + k - 1$. От тези позиции трябва да изберем $k-1$ места, на които да поставим преградите (или еквивалентно, да изберем $n$ места за звездите). 
Следователно, броят на решенията е:
\[ \binom{n+k-1}{k-1} = \binom{n+k-1}{n} \]

Този резултат е изключително важен и се прилага в много контексти, включително за намиране на броя на разпределенията на $n$ неразличими частици в $k$ различими клетки (където редът на частиците в дадена клетка няма значение).

\section{Просто случайно блуждаене и Принцип на отражението}

Ще разгледаме един класически модел в теорията на вероятностите - просто случайно блуждаене (Simple Random Walk).
Нека една частица се движи по целочислените точки на реалната права. Тя стартира от началото на координатната система (точка 0) в момент $t=0$. Във всеки следващ дискретен момент от времето $t=1, 2, \dots$ частицата прави стъпка надясно (към $+1$) или наляво (към $-1$).
Нека $\xi_i \in \{+1, -1\}$ е стойността на $i$-тата стъпка. Тогава позицията на частицата след $n$ стъпки е:
\[ S_n = \xi_1 + \xi_2 + \dots + \xi_n \]
като по дефиниция приемаме $S_0 = 0$.

Можем да представим блуждаенето графично в координатната равнина $(t, S_t)$, където по хоризонталната ос нанасяме времето (броя стъпки), а по вертикалната - позицията на частицата. Тогава блуждаенето представлява начупена линия, започваща от $(0,0)$. Всяка отсечка от тази линия свързва точките $(i-1, S_{i-1})$ и $(i, S_i)$.

\subsection{Брой на всички пътища}
Ако разгледаме блуждаене с дължина $n$, то броят на всички възможни пътища (траектории) е $2^n$, тъй като на всяка от $n$-те стъпки имаме точно 2 възможни избора ($+1$ или $-1$).

За да стигнем от точка $(0,0)$ до точка $(n, x)$, трябва броят на стъпките надясно ($+1$) и наляво ($-1$) да са съответно $p$ и $q$, такива че да са изпълнени следните две условия:
\begin{align*}
p + q &= n \\
p - q &= x
\end{align*}
Като съберем двете уравнения, намираме:
\[ 2p = n + x \implies p = \frac{n+x}{2} \]
За да съществува такъв път, е необходимо и достатъчно $n+x$ да бъде четно число и $-n \leq x \leq n$.
В този случай, броят на пътищата от $(0,0)$ до $(n,x)$ е равен на броя на начините да изберем $p$ стъпки надясно от общо $n$ налични стъпки, тоест:
\[ N_n(x) = \binom{n}{p} = \binom{n}{\frac{n+x}{2}} \]

\subsection{Принцип на отражението}
Често ни интересуват пътища, които удовлетворяват определени условия, например пътища, които остават неотрицателни ($S_i \geq 0$ за всяко $i$). За да преброим такива пътища, използваме Принципа на отражението, въведен от Дезире Андре.

Нека разгледаме множеството от всички пътища, започващи от точка $(0,x)$ до $(n,y)$, където $x, y > 0$. Искаме да намерим броя на пътищата, които докосват или пресичат хоризонталната ос (т.е. $S_i = 0$ за някое $i$). Нека ги наречем "лоши" пътища.

Вземаме един произволен такъв лош път. Нека $t_0$ е първият момент, в който пътят докосва нулата (т.е. $S_{t_0} = 0$, но $S_i > 0$ за $i < t_0$). Принципът на отражението гласи, че ако отразим симетрично началната част от пътя (от $t=0$ до $t=t_0$) спрямо абсцисната ос $S=0$, ще получим нов път, който започва от $(0,-x)$, докосва нулата в $t_0$ и продължава по оригиналния път до $(n,y)$. 

Тъй като всеки път от $(0,-x)$ до $(n,y)$ (при $x,y > 0$) задължително пресича оста $S=0$ в някакъв момент, съществува взаимно еднозначно съответствие (биекция) между лошите пътища от $(0,x)$ до $(n,y)$ и множеството на всички възможни пътища от $(0,-x)$ до $(n,y)$.

Следователно, броят на лошите пътища от $(0,x)$ до $(n,y)$ е точно равен на общия брой пътища от $(0,-x)$ до $(n,y)$.

\subsection{Пътища, които остават неотрицателни}
Нека разгледаме пътищата, стартиращи от $(0,0)$ и завършващи в $(2n, 0)$. Общият им брой, както вече изведохме, е $N_{2n}(0) = \binom{2n}{n}$.

Интересуваме се от броя на пътищата от $(0,0)$ до $(2n, 0)$, които остават неотрицателни през цялото време, т.е. $S_i \geq 0$ за $i=1, \dots, 2n$. 
За да намерим техния брой, ще преброим "лошите" пътища – тези, за които $S_i = -1$ за поне едно $i$ – и ще ги извадим от общия брой пътища.

За всеки лош път от $(0,0)$ до $(2n,0)$, съществува първи момент $t_0$, в който пътят достига ниво $-1$. Според принципа на отражението, ако отразим пътя след момент $t_0$ спрямо правата $S = -1$, ще получим път, който започва от $(0,0)$ и завършва в $(2n, -2)$. 
Още по-удобно, ако транслираме или използваме биекция от първата стъпка: лошите пътища са в биекция с всички пътища от $(0,-2)$ до $(2n,0)$.
Броят на пътищата от $(0,-2)$ до $(2n,0)$ се намира като решим:
\begin{align*}
p + q &= 2n \\
p - q &= 2
\end{align*}
Оттук $2p = 2n + 2 \implies p = n+1$. Следователно броят им е $\binom{2n}{n+1}$.

Тогава броят на "добрите" пътища (които никога не достигат $-1$, следователно $S_i \geq 0$ за всички $i$) от $(0,0)$ до $(2n,0)$ е:
\[ \binom{2n}{n} - \binom{2n}{n+1} \]
Тази разлика може да се опрости алгебрично по следния начин:
\begin{align*}
\binom{2n}{n} - \binom{2n}{n+1} &= \frac{(2n)!}{n!n!} - \frac{(2n)!}{(n+1)!(n-1)!} \\
&= \frac{(2n)!}{n!(n-1)!} \left( \frac{1}{n} - \frac{1}{n+1} \right) \\
&= \frac{(2n)!}{n!(n-1)!} \left( \frac{n+1 - n}{n(n+1)} \right) \\
&= \frac{(2n)!}{n!(n-1)!} \frac{1}{n(n+1)} \\
&= \frac{1}{n+1} \frac{(2n)!}{n!n!} \\
&= \frac{1}{n+1} \binom{2n}{n}
\end{align*}
Този забележителен резултат представлява $n$-тото число на Каталан ($C_n$). Числата на Каталан имат огромно значение в комбинаториката, тъй като те броят множество различни комбинаторни структури (например правилни скобови изрази, триангулации на изпъкнали многоъгълници и дървета).

\subsection{Пътища с условие за строга положителност}
Ако разгледаме по-строго условие – да намерим броя на пътищата от $(0,0)$ до $(2n,0)$, които се връщат в нулата за първи път чак на последната стъпка $2n$ (т.е. $S_i > 0$ за $1 \leq i \leq 2n-1$ и $S_{2n} = 0$).

Тогава първата стъпка на такъв път задължително е нагоре към $(1,1)$, а предпоследната позиция задължително е $(2n-1, 1)$, откъдето прави стъпка надолу към $(2n,0)$.
Броят на тези пътища съвпада с броя на пътищата от $(1,1)$ до $(2n-1, 1)$, които никога не докосват нулата (остават строго положителни, $S_i > 0$). 

Ако транслираме координатната система с вектор $(-1,-1)$, това е равносилно на намирането на броя на пътищата от $(0,0)$ до $(2n-2, 0)$, които остават неотрицателни (т.е. не слизат под новото ниво 0).
Според предходната точка, този брой е равен на съответното число на Каталан за $2n-2$ стъпки:
\[ \frac{1}{n} \binom{2n-2}{n-1} = C_{n-1} \]
което е $(n-1)$-тото число на Каталан.

\end{document}